\documentclass[11pt,a4paper]{article}
\usepackage[margin=22mm]{geometry}
\usepackage[T1]{fontenc}
\usepackage[utf8]{inputenc}
\usepackage{lmodern,amsmath,booktabs,longtable,array,graphicx,microtype,xcolor}
\usepackage[hidelinks,breaklinks]{hyperref}
\usepackage{url}
\setlength{\emergencystretch}{3em}
\title{Deadline Forecasting Under a Fractional CPU Quota: Prospective Comparisons and Calibration-Transfer Limits}
\author{}
\date{}
\begin{document}
\maketitle
\begin{abstract}

Accurate average capacity need not preserve deadline forecasts, and nominal
calibration controls need not assure later response-distribution transfer.
We evaluate established approaches in six separately calibrated prospective
campaigns on one Linux host under a 25 ms/100 ms CPU quota. In a Go HTTP
service, eighteen steady and burst trials give 25 ms deadline-miss mean absolute
errors of 73.599 percentage points for the tested constant-rate busy-capacity
FIFO and 2.341 for periodic service with calibrated delay. A separate eighteen-trial
steady-traffic study gives 10 ms error of 1.844 points for a queue-adjusted CDF,
versus 67.340 for periodic+delay, in an observed low-queue domain. A retrospective
complete-response CDF control gives 1.802 points, limiting decomposition claims;
the empirical reference still fails on exploratory burst traffic. In stock
NGINX gzip compression, a complete-response reference at the frozen 20.565395 ms
deadline gives 5.981 points error, with conditional 95\% interval [1.321,10.642],
failing the required upper bound <=5 points. Fresh controls retain this adverse
result: matching assigned power policy reduces error from 24.729 to 15.708
points but remains insufficiently accurate. Later fixed-policy, rate-matched
calibration gives 1.194 points error, while ordinary 5/s calibration also passes
at 1.972 points; the 0.778-point improvement interval [-0.375,1.931] does not
support the required advantage. A fresh five-deadline panel gives policy-specific
worst-cutoff block scores of 6.611 and 4.889 points, with conditional 97.5\%
interval upper bounds 11.391 and 7.158; neither meets the 5-point accuracy gate.
Zero error at an older saturated deadline therefore does not establish panel
accuracy. Retrospective reliability screens supply no confirmed advance selector;
measurement qualification and precision/cost limits hold further local collection.
The contribution is reproducible prospective
evidence of the tested procedures' adequacy and transfer limitations. The
campaigns establish no new estimator, historical physical cause, universal
selection rule or independently verified hardware generality.

\end{abstract}
\section{Question and contribution}

We ask: \textbf{how accurately do constant-rate, periodic-budget and empirical latency references forecast deadline misses under the tested fractional CPU quota, deadlines and traffic conditions, and what do ordinary calibration-transfer controls establish?}

The completed capacity-calibration paper \cite{capacityprecursor} compared independently measured
capacity inputs within a common finite FIFO model. This extension tests the
next assumption: whether replacing intermittently available service with its
average processing rate preserves a useful deadline forecast. It also asks
what happens when a deadline approaches the measured response-time
distribution itself.

The artifact labels V1/V29 denote precursor capacity studies, V30 the 25 ms
comparison, and V31 the subsequent 10 ms study. We retain these identifiers
to make each claim traceable to its separate frozen protocol.

There are eight distinct evidence stages:

\begin{enumerate}
\item The retained V1 and V29 studies establish scoped capacity and long-deadline
evidence with independently calibrated inputs.
\item V30 prospectively confirms a particular 25 ms prediction/decision contrast
on fresh traffic patterns, after the ordinary timing approximation was
developed using V29 discovery data.
\item The new empirical latency-CDF comparison is retrospective development
prompted by V30's 10 ms failure. Its parameters use calibration records,
but its method was chosen after validation outcomes; it is not another
prospective result.
\item V31 prospectively confirms the predeclared queue-adjusted calibration CDF at 10 ms
using fresh calibration and eighteen trials at new steady rates. The V30
exploratory evidence remains discovery; V31 has its own sealed predictions,
technical qualification, domain checks and confirmation gates.
\item A separate stock-NGINX campaign tests a newly scoped complete-response CDF
forecast and does not meet its primary accuracy criterion. It shares the
physical host and broad client-completion outcome, with different calibration
inputs, completion checks, observed domain and model comparison.
\item After a stopped exploratory power-policy pilot, a fresh all-offer confirmation
compares pooled and policy-matched references. Matching improves error but
fails the frozen absolute-accuracy requirement; the pilot is not pooled in.
\item A fresh fixed-P0 study compares target-rate and ordinary 5/s references.
Both transfer accurately, without the required incremental advantage from
rate matching. This later episode does not identify or repair the earlier cause.
\item A fresh policy-specific five-deadline panel tests the scope suggested by
retrospective deadline checks. Neither policy meets its frozen accuracy gate,
despite substantial matching advantage over pooling and zero P1 error at the
older saturated cutoff. It is a separate prospective study on the same host.
\end{enumerate}
This separation prevents a successful exploratory remedy from acquiring
confirmation status merely because its parameters came from calibration.

The contribution has three parts: prospective numerical comparisons of selected
approximations at fixed deadlines and declared traffic/domain contracts; a
calibration-transfer audit retaining the adverse NGINX result and complete
ordinary-control follow-ups; and an auditable artifact preserving fresh
calibration, advance forecast barriers, all-offer outcomes and domain exits.
The artifact strengthens reproducibility, not algorithmic novelty. These
studies quantify the adequacy of selected existing approaches. The
Go and NGINX validation both count all-offer client-completion misses, including
waiting. Different completion checks, calibration inputs, deadlines, domains
and comparison contracts prevent combining their results into a controlled
application effect or a common accuracy interval.

\begin{figure}[htbp]\centering
{\footnotesize\setlength{\unitlength}{1mm}\begin{picture}(160,87)
{\color{black}\linethickness{.2mm}\qbezier(8.000,22.000)(27.000,22.000)(46.000,22.000)}
{\color{black}\linethickness{.2mm}\qbezier(8.000,22.000)(8.000,46.000)(8.000,70.000)}
{\color{gray!25}\linethickness{.1mm}\qbezier(8.000,22.000)(27.000,22.000)(46.000,22.000)}
\put(6.500,22.000){\makebox(0,0)[r]{0}}
{\color{gray!25}\linethickness{.1mm}\qbezier(8.000,33.707)(27.000,33.707)(46.000,33.707)}
\put(6.500,33.707){\makebox(0,0)[r]{20}}
{\color{gray!25}\linethickness{.1mm}\qbezier(8.000,45.415)(27.000,45.415)(46.000,45.415)}
\put(6.500,45.415){\makebox(0,0)[r]{40}}
{\color{gray!25}\linethickness{.1mm}\qbezier(8.000,57.122)(27.000,57.122)(46.000,57.122)}
\put(6.500,57.122){\makebox(0,0)[r]{60}}
{\color{gray!25}\linethickness{.1mm}\qbezier(8.000,68.829)(27.000,68.829)(46.000,68.829)}
\put(6.500,68.829){\makebox(0,0)[r]{80}}
\put(13.000,22){\color{gray}\rule{9mm}{43.0824mm}}
\put(17.500,67.082){\makebox(0,0)[c]{73.599}}
\put(31.000,22){\color{blue!65!black}\rule{9mm}{1.3706mm}}
\put(35.500,25.371){\makebox(0,0)[c]{2.341}}
\put(27.000,79.000){\makebox(0,0)[c]{V30: Go, 25 ms}}
\put(17.500,17.000){\makebox(0,0)[c]{Constant-rate}}
\put(17.500,13.000){\makebox(0,0)[c]{FIFO}}
\put(35.500,17.000){\makebox(0,0)[c]{Periodic}}
\put(35.500,13.000){\makebox(0,0)[c]{+ delay}}
{\color{black}\linethickness{.2mm}\qbezier(62.000,22.000)(81.000,22.000)(100.000,22.000)}
{\color{black}\linethickness{.2mm}\qbezier(62.000,22.000)(62.000,46.000)(62.000,70.000)}
{\color{gray!25}\linethickness{.1mm}\qbezier(62.000,22.000)(81.000,22.000)(100.000,22.000)}
\put(60.500,22.000){\makebox(0,0)[r]{0}}
{\color{gray!25}\linethickness{.1mm}\qbezier(62.000,33.707)(81.000,33.707)(100.000,33.707)}
\put(60.500,33.707){\makebox(0,0)[r]{20}}
{\color{gray!25}\linethickness{.1mm}\qbezier(62.000,45.415)(81.000,45.415)(100.000,45.415)}
\put(60.500,45.415){\makebox(0,0)[r]{40}}
{\color{gray!25}\linethickness{.1mm}\qbezier(62.000,57.122)(81.000,57.122)(100.000,57.122)}
\put(60.500,57.122){\makebox(0,0)[r]{60}}
{\color{gray!25}\linethickness{.1mm}\qbezier(62.000,68.829)(81.000,68.829)(100.000,68.829)}
\put(60.500,68.829){\makebox(0,0)[r]{80}}
\put(67.000,22){\color{blue!65!black}\rule{9mm}{39.4183mm}}
\put(71.500,63.418){\makebox(0,0)[c]{67.340}}
\put(85.000,22){\color{orange!85!black}\rule{9mm}{1.0797mm}}
\put(89.500,25.080){\makebox(0,0)[c]{1.844}}
\put(81.000,79.000){\makebox(0,0)[c]{V31: Go, 10 ms}}
\put(71.500,17.000){\makebox(0,0)[c]{Periodic}}
\put(71.500,13.000){\makebox(0,0)[c]{+ delay}}
\put(89.500,17.000){\makebox(0,0)[c]{Queue-adj.}}
\put(89.500,13.000){\makebox(0,0)[c]{CDF}}
{\color{black}\linethickness{.2mm}\qbezier(116.000,22.000)(135.000,22.000)(154.000,22.000)}
{\color{black}\linethickness{.2mm}\qbezier(116.000,22.000)(116.000,46.000)(116.000,70.000)}
{\color{gray!25}\linethickness{.1mm}\qbezier(116.000,22.000)(135.000,22.000)(154.000,22.000)}
\put(114.500,22.000){\makebox(0,0)[r]{0}}
{\color{gray!25}\linethickness{.1mm}\qbezier(116.000,37.211)(135.000,37.211)(154.000,37.211)}
\put(114.500,37.211){\makebox(0,0)[r]{5}}
{\color{gray!25}\linethickness{.1mm}\qbezier(116.000,52.423)(135.000,52.423)(154.000,52.423)}
\put(114.500,52.423){\makebox(0,0)[r]{10}}
{\color{gray!25}\linethickness{.1mm}\qbezier(116.000,67.634)(135.000,67.634)(154.000,67.634)}
\put(114.500,67.634){\makebox(0,0)[r]{15}}
\put(125.000,32.479){\color{black}\circle*{.9}}
\put(127.000,66.958){\color{black}\circle*{.9}}
\put(129.000,33.831){\color{black}\circle*{.9}}
\put(131.000,32.817){\color{black}\circle*{.9}}
\put(133.000,41.268){\color{black}\circle*{.9}}
\put(135.000,33.831){\color{black}\circle*{.9}}
{\color{orange!85!black}\linethickness{.4mm}\qbezier(130.000,26.019)(130.000,40.197)(130.000,54.376)}
{\color{orange!85!black}\linethickness{.4mm}\qbezier(128.000,26.019)(130.000,26.019)(132.000,26.019)}
\put(135.000,26.019){\makebox(0,0)[l]{1.321}}
{\color{orange!85!black}\linethickness{.4mm}\qbezier(128.000,54.376)(130.000,54.376)(132.000,54.376)}
\put(135.000,54.376){\makebox(0,0)[l]{10.642}}
\put(130.000,40.197){\color{orange!85!black}\circle*{1.7}}
\put(135.000,40.197){\makebox(0,0)[l]{5.981}}
{\color{red!65!black}\linethickness{.2mm}\qbezier(116.000,37.211)(117.000,37.211)(118.000,37.211)}
{\color{red!65!black}\linethickness{.2mm}\qbezier(120.000,37.211)(121.000,37.211)(122.000,37.211)}
{\color{red!65!black}\linethickness{.2mm}\qbezier(124.000,37.211)(125.000,37.211)(126.000,37.211)}
{\color{red!65!black}\linethickness{.2mm}\qbezier(128.000,37.211)(129.000,37.211)(130.000,37.211)}
{\color{red!65!black}\linethickness{.2mm}\qbezier(132.000,37.211)(133.000,37.211)(134.000,37.211)}
{\color{red!65!black}\linethickness{.2mm}\qbezier(136.000,37.211)(137.000,37.211)(138.000,37.211)}
{\color{red!65!black}\linethickness{.2mm}\qbezier(140.000,37.211)(141.000,37.211)(142.000,37.211)}
{\color{red!65!black}\linethickness{.2mm}\qbezier(144.000,37.211)(145.000,37.211)(146.000,37.211)}
{\color{red!65!black}\linethickness{.2mm}\qbezier(148.000,37.211)(149.000,37.211)(150.000,37.211)}
{\color{red!65!black}\linethickness{.2mm}\qbezier(152.000,37.211)(153.000,37.211)(154.000,37.211)}
\put(135.000,79.000){\makebox(0,0)[c]{NGINX: 20.565395 ms}}
\put(135.000,17.000){\makebox(0,0)[c]{Full-response CDF}}
\put(135.000,13.000){\makebox(0,0)[c]{Conditional 95\% interval}}
\put(80.000,4.000){\makebox(0,0)[c]{Mean absolute miss-fraction error (percentage points); independent panel scales}}
\end{picture}}
\caption{The initial three separate prospective comparisons. Bars show equal-profile/block mean absolute miss-fraction error. The NGINX point and interval show the declared conditional six-block mean-error interval; small points are its six block errors. The dashed line is its 5 pp upper-interval criterion. Completion checks, calibration inputs, deadlines, domains, comparators and criteria differ; these panels do not establish a common model ranking. The retrospective complete-response control is reported in Section 10; the fresh policy and rate controls are reported in Sections 12 and 13.}\label{fig:primary}
\end{figure}

\section{Related work and contribution scope}

Periodic-server response-distribution modeling predates this extension.
\href{https://www.cse.wustl.edu/~cdgill/publications/rtss19.pdf}{Li, Lu and Gill (RTSS 2019)}
analyze Poisson arrivals with deterministic demand under a deferrable server
and evaluate synthetic and Redis services on Xen. These premises differ from
our finite schedules and ordinary Linux bandwidth ceiling. Our periodic model
is a simplified baseline, not their implemented analysis. This review inspected
indexed primary passages; direct full-PDF retrieval was unavailable.

The average-speed approximation in
\href{https://retis.santannapisa.it/people/nino/publication/TCC23.pdf}{Cucinotta et al. (TCC 2023)}
explicitly assumes a reservation period small relative to request processing
time, under HCBS temporal isolation. Its evaluation uses a 2 ms period. Our
100 ms ceiling and short requests do not reproduce this regime, and failure
of our constant-rate comparator does not refute their method.

Empirical profiling for tail prediction is also established.
\href{https://www2.cs.uh.edu/~fyan/Paper/Feng-ICAC18.pdf}{CEDULE (ICAC 2018)}
fits quantile models using CPU-limit/load profiles for burstable-cloud
scheduling. Its richer conditioning and quantile metrics differ from our
fixed empirical miss-probability reference. We have neither implemented a
matched comparison nor demonstrated superiority to that method.

\href{https://link.springer.com/content/pdf/10.1007/s13174-010-0012-9.pdf}{Mi et al. (JISA 2010)}
study temporal dependence and limitations of average models, fitting richer
service models from busy-period observations. The broad importance of
burstiness and dependence is therefore background, not our discovery.
Likewise, \href{https://arxiv.org/html/2407.11582v1}{Jeffery et al. (2024 preprint)}
already distinguish queue and worker latency in quota-aware thread management;
stage instrumentation itself supplies no algorithmic novelty.

\href{https://syslab.cs.washington.edu/papers/latency-socc14.pdf}{Tales of the Tail (SoCC 2014)} \cite{tales}
already measures hardware, OS and application sources of tail latency, including
NGINX and low-utilization power effects. Our assigned-policy comparison does
not discover that power management affects tails or identify a new mechanism.
\href{https://people.csail.mit.edu/delimitrou/papers/2022.hpca.retail.pdf}{ReTail (HPCA 2022)} \cite{retail}
uses request features and frequency-conditioned latency models with active
frequency control, monitoring and drift-triggered retraining. Those richer
inputs and online decisions differ from our campaign-wide frozen CDF.
\href{https://drops.dagstuhl.de/storage/00lipics/lipics-vol375-ecrts2026/html/LIPIcs.ECRTS.2026.17/LIPIcs.ECRTS.2026.17.html}{MPORA (ECRTS 2026)} \cite{mpora}
uses execution-state/allocation models and weighted conformal uncertainty for
deadline-aware shared-resource allocation. Its task-progress endpoint and
profiling contract differ; shifted coverage remains assumption-dependent.
Neither our CDF failures nor our lack of those inputs refute these methods.

The possible incremental contribution is the measured, prospectively tested
contrast across selected deadlines and traffic conditions, accompanied by
adverse transfer evidence and fresh ordinary-control follow-ups. It is restricted to the declared comparators and
observed domains. The queue-adjusted CDF does not dominate every deadline,
a steep CDF is not a new mechanism, and NGINX tests a separate calibration-transfer
contract for the broad client-response outcome.
No universal deadline/period boundary, advance applicability selector,
independent-hardware replication or state-of-the-art superiority follows.
The original focused review and the later transfer-consolidation capability
audit retain source locators and access/coverage gaps in the supporting material.
Finding no exact match in selected passages is insufficient to claim a first
result; scientific novelty remains a candidate requiring further assessment.

\section{Model and prior-method boundary}

The constant-rate FIFO uses calibrated throughput $C$, assigning each admitted
job service time $1/C$. The periodic-budget approximation consumes an effective
demand $d = (B/P)/C$ while credit is available, and replenishes its global budget
$B$ at each period boundary $P$. Here $B$ denotes CPU credit and $Q$ below denotes
measured request queue wait. Both model one job in service and up to 128 waiting
jobs; modeled rejection counts as a deadline miss.

Periodic CPU supply and average-speed approximations are established ideas.
\href{https://docs.kernel.org/scheduler/sched-bwc.html}{Linux bandwidth-control documentation}
describes quota ceilings, replenishment, allocation slices and retained
CPU-local runtime. Our single global-budget abstraction omits those local
details, competing service threads and execution-cost variation.
\href{https://retis.santannapisa.it/people/nino/publication/TCC23.pdf}{Cucinotta et al.}
study hierarchical real-time reservations, with guarantees and modeling
premises that differ from this ordinary bandwidth-limited benchmark. Our
experiment is not a reproduction or refutation of their full method.

The main periodic forecast averages 100 shared phases at 1 ms spacing, with
full initial credit. This is an assumed phase prior; kernel replenishment phase
was neither observed nor aligned. A response-delay distribution is fitted
from four low-rate calibration trials using e2e minus queue minus worker-wall
time, restricted to successful calibration requests with queue wait <=100 us.
It is appended independently of modeled service. Neither effective demand nor
the residual is an isolated measurement of SHA CPU demand.

\section{Evidence sequence and experimental contract}

\subsection{Outcome and statistical units}

For all five primary campaigns, an offer at deadline D misses if
it is unsuccessful or its successful
scheduled-offer-to-completion duration exceeds D; equality is timely. Failed,
rejected, timed-out or missing offers retain their place in the denominator.
Let $p_{b,k}(D)$ be the observed all-offer miss fraction for profile $k$ in block $b$
and $\hat p_{b,k}(D)$ its frozen forecast. Trial absolute error is
$100\left|\hat p_{b,k}(D)-p_{b,k}(D)\right|$ percentage points. Each block error is the
arithmetic mean of its three profile errors, and primary MAE is the mean of
the six block errors in the initial studies. The power confirmation averages
four policy/rate cells per block; rate transfer averages two rate cells. Equal
cell weighting differs from pooling all offers.

For V30 and V31, improvement is the difference between comparator and tested
reference block errors. The declared 95\% interval is
$\bar\Delta \pm 2.571\,s_\Delta/\sqrt{6}$, where $s_\Delta$ is the sample
standard deviation of the six block improvements. NGINX applies the same interval
formula to its six block errors rather than an improvement contrast. These
model-based intervals assume approximately independent, normally distributed
block summaries and condition on the finite frozen calibrations. They do not
propagate calibration uncertainty or treat requests as independent replicates.
Blocks were collected sequentially. Counterbalancing controls profile position
but cannot establish independence or remove drift between blocks.

The two fresh NGINX controls apply the same conditional interval formula to
their own block errors and paired improvements. Matched accuracy requires an
error-interval upper bound <=5 pp. Incremental advantage requires mean
improvement >=2 pp and an improvement-interval lower bound >0; it does not
require a >=2 pp lower confidence bound. Accuracy and advantage are reported
separately. These are frozen engineering references, not universal guarantees.
Engineering tolerances and the six-block budget were fixed before outcomes;
they are not established universal SLOs or power guarantees.

\subsection{Resource, calibration and validation contracts}

V29 comprises four conditions of the same service and host: CPU allowance
25\%/50\% crossed with two SHA iteration counts. Each complete condition uses
sixteen calibration and sixteen validation trials. Three complete conditions
come from the overnight generation; the final condition comes from an
independently recalibrated rerun. The interrupted final-condition record is
preserved as supplementary evidence and is not pooled into the full matrix.
Conditions ran sequentially, so cross-condition differences are confounded
with time. Request-size changes do not constitute independent application
transfer.

V30 uses $B=25$ ms, $P=100$ ms, zero configured burst credit, 88000 SHA iterations,
one worker, GOMAXPROCS=1, queue size 128 and server CPUs 2 and 4. Client/control
work uses CPUs 0 and 1. Client concurrency is 512 and timeout is 5 s. Six
functional cases and an unreleased resource preflight qualify the retained
binaries before scientific collection. Fresh observer inventory, clocks,
power/thermal settling, ownership and the 1 Hz host monitor are retained.

The earlier host record identifies an AMD Ryzen 5 5600H with twelve logical
CPUs and approximately 16 GB RAM. That hardware description comes from the
retained precursor record, rather than a new hardware inventory in each
campaign. V30 and V31's own observer audits record Linux 7.2.8-arch1-2, and
their service records and client-build manifests identify stock Go 1.26.1.
Both campaigns record \path{amd-pstate-epp}, the \path{powersave} governor and
\path{balance_performance} preference; CPU frequency is not fixed. The request
initializes a 32-byte state to bytes 0 through 31 and applies SHA-256 to the
preceding 32-byte state 88,000 times. Each scientific slot uses a fresh service,
with instrumented queue/completion records and runtime tracing disabled.

Both campaigns independently collect the same sixteen-slot calibration design:
two 30 s idle trials, two 30 s steady trials at each of 3.75, 7.5, 15, 22.5 and
30 requests/s, and four 60 s mixed trials. Each mixed trial has four 15 s cycles,
with 10 s at 3.75 or 7.5 requests/s and 5 s at 42.45 requests/s; each low rate
occurs twice. Warmup is 20 s, followed by drain and metrics reset. Busy capacity
uses completion counts in [11.5,14.5) s and the corresponding windows shifted
by 15, 30 and 45 s. Clock-expanded windows must have a strictly positive
waiting queue throughout; the capacity estimate is the arithmetic mean of the
four trial rates. The declared maximum/minimum capacity spread is <=5\%.
Steady qualification uses windows [3,10), [10,17) and [17,24) s, completion-rate
bounds within 5\% of the offered rate, peak waiting queue <=1 at 3.75/7.5
requests/s and <=5 at higher rates, and zero measurement throttled periods
at the two low rates. The four low-rate cohorts, each measured for 30 s,
supply the timing fit after selecting successful queue<=100 us requests.
These contracts are inherited from the sealed capacity design; fresh V30 and
V31 measurements are not pooled.

Sixteen fresh calibration trials precede a forecast seal. Eighteen validation
trials use six blocks in order ABC, BCA, CAB, ACB, CBA, BAC. Each profile occurs
in each position twice. Measurement is 60 s after a 20 s warmup, with reset and
bounded drain between trials.

{\small\setlength{\tabcolsep}{2pt}\begin{longtable}{>{\raggedright\arraybackslash}p{0.25200\linewidth}>{\raggedright\arraybackslash}p{0.32400\linewidth}>{\raggedright\arraybackslash}p{0.32400\linewidth}}\toprule
Profile & Traffic & Offers per trial \\
\midrule\endhead
A & Steady 4 requests/s & 240 \\
B & Steady 18 requests/s & 1080 \\
C & Repeated 1 s at 4 requests/s, then 1 s at 42 requests/s & 1380 \\
\bottomrule\end{longtable}}

There are 16200 validation offers. The endpoint measures scheduled offer to
successful client completion; failed or missing offers count as misses.
The Go client timestamps completion after consuming the HTTP response body
and closing it. NGINX timestamps complete encoded-body receipt before closing,
decoding and hash validation; correctness failures still count as unsuccessful
offers. Both validation outcomes include queue waiting. V31 never subtracts
the validation request's queue interval from its deadline outcome.
Empty queues are legitimate validation behavior. Six blocks, rather than
requests, are the replication units. Counterbalancing does not prove their
independence.

The primary deadline is 25 ms and the experimental miss budget is 5\%. Frozen
confirmation gates require all eighteen qualified trials, <=3 pp mean block
MAE for periodic+delay, positive improvement in every block and >=5 pp at the
lower endpoint of the declared 95\% block-based t interval, correct miss-budget
classification for every trial, and at least one profile meeting the target in
every block. The interval uses df=5 and multiplier 2.571, assuming approximately
independent normal block contrasts. These are declared engineering targets;
the trial count is a feasibility budget rather than a demonstrated power
calculation.

\section{Retained capacity and long-deadline evidence}

{\small\setlength{\tabcolsep}{2pt}\begin{longtable}{>{\raggedright\arraybackslash}p{0.25200\linewidth}>{\raggedright\arraybackslash}p{0.32400\linewidth}>{\raggedright\arraybackslash}p{0.32400\linewidth}}\toprule
V29 condition & Busy capacity requests/s & Busy-model 500 ms MAE pp \\
\midrule\endhead
50\% CPU, 44000 iterations & 153.792 & 0.479 \\
25\% CPU, 44000 iterations & 69.625 & 1.359 \\
50\% CPU, 88000 iterations & 77.646 & 0.672 \\
25\% CPU, 88000 iterations & 34.854 & 1.295 \\
\bottomrule\end{longtable}}

The frozen busy-comparator gates pass in all four complete V29 conditions.
This extends within-service robustness evidence but neither establishes
causal cross-condition effects nor overturns V1's inconclusive minimum-effect
comparison against its quadratic benchmark. Each study retains its own
calibration, traffic and gates.

\section{Prospective 25 ms result}

V30's fresh busy estimate is 34.917 requests/s. The constant-rate substitution
therefore assigns every job 28.640 ms of service and predicts 100\% misses at
25 ms even without a queue. Actual steady requests frequently finish sooner
than that imposed floor.

The comparator's 100\% forecast follows directly from its deterministic
service floor; this is a deliberately simple substitution, not a competitive
implementation of a published tail model. The practical finding is that this
busy-capacity input can be adequate for the retained long-deadline forecasts
yet impose an inappropriate short-deadline floor when treated as continuous
speed. The quantified contrast is a case-specific calibration/model mismatch,
not a new scheduling law or evidence that published periodic analyses fail.

{\small\setlength{\tabcolsep}{2pt}\begin{longtable}{>{\raggedright\arraybackslash}p{0.25200\linewidth}>{\raggedright\arraybackslash}p{0.21600\linewidth}>{\raggedright\arraybackslash}p{0.21600\linewidth}>{\raggedright\arraybackslash}p{0.21600\linewidth}}\toprule
Traffic & Observed mean miss fraction & Periodic+delay forecast & Periodic+delay MAE pp \\
\midrule\endhead
Steady 4 requests/s & 0.000\% & 0.000\% & 0.000 \\
Steady 18 requests/s & 0.000\% & 0.000\% & 0.000 \\
Bursty profile C & 79.203\% & 72.179\% & 7.024 \\
\bottomrule\end{longtable}}

The primary equal-profile/block MAE is 73.599 pp for fluid busy and 2.341 pp
for periodic+delay. Every block improves; mean paired improvement is 71.258 pp
with declared interval [68.900,73.615] pp. All eighteen observed <=5\%/>5\%
decisions match the periodic forecast. All 34 scientific trials complete
without a technical retry; independent completion replay passes.

The primary accuracy gate applies to the mean. Two blocks exceed 3 pp and the
maximum individual error is 12.097 pp. Burst-only MAE is 7.024 pp; weighting
by offered requests instead of profiles gives 3.590 pp. Thus the result is
useful classification of widely separated operating conditions, not evidence
of <=3 pp error on every traffic pattern or a precise boundary near 5\% misses.
Zero observed misses in steady trials do not guarantee zero future misses.

\section{Deadline-resolution limits}

{\small\setlength{\tabcolsep}{2pt}\begin{longtable}{>{\raggedright\arraybackslash}p{0.25200\linewidth}>{\raggedright\arraybackslash}p{0.32400\linewidth}>{\raggedright\arraybackslash}p{0.32400\linewidth}}\toprule
Deadline ms & Fluid busy MAE pp & Periodic+delay MAE pp \\
\midrule\endhead
10 & 24.394 & 49.622 \\
20 & 73.080 & 2.318 \\
25 (primary) & 73.599 & 2.341 \\
35 & 4.018 & 3.306 \\
50 & 2.947 & 4.122 \\
75 & 1.385 & 6.181 \\
100 & 0.918 & 8.173 \\
150 & 2.738 & 9.691 \\
500 & 0.000 & 0.000 \\
\bottomrule\end{longtable}}

The periodic approximation improves 20\textendash{}25 ms forecasts, but fails at 10 ms and
is less accurate than fluid busy at 50\textendash{}150 ms. The 500 ms endpoint has no misses
and supplies no comparative discrimination in V30. Secondary endpoints carry
no new superiority-significance claim.

The fixed zero-phase/full-credit diagnostic has 0.902 pp MAE at 25 ms, below the
main model's 2.341 pp. This was a prespecified scenario, not measured phase;
promoting it after outcomes would require a new selection/confirmation
contract. Periodic without response delay has 2.729 pp MAE, so this experiment
does not demonstrate that delay calibration is necessary. The discovery
model-switch rule creates nonmonotone miss curves in six trials. It is not one
coherent latency distribution and cannot be presented as a general correction.

\section{Exploratory timing comparison at 10 ms}

All twelve steady validation trials record zero throttled periods. Median queue
wait is about 0.009 ms. Nevertheless, about 67\% of their requests miss 10 ms.
This failure appears without recorded quota throttling or substantial queues.

The model infers 7.160 ms effective CPU demand from busy throughput. The four
low-rate calibration cohorts instead have 8.855 ms median worker-wall time,
1.397 ms mean non-queue residual and 10.228 ms median paired queue-free
latency. The deadline lies near that distribution's center; an error of one or
two milliseconds materially changes the predicted miss fraction.

For each selected calibration request define $R=E-Q$, where $E$ is complete
client response time and $Q$ is observed queue wait. If $W$ is worker-wall time,
the retained worker/residual sum is exactly $W+(E-Q-W)=E-Q$. The primary CDF
does not require the worker/residual split; that split serves stage diagnostics
and the independent-marginal convolution control. The ordinary empirical
estimate is $\hat p_R(D)=n^{-1}\sum_{i=1}^{n}\mathbf{1}\{R_i>D\}$.
There are 676 samples across four independently collected calibration trials.
No validation residual, phase or threshold is fitted. The choice to examine
this model is post-V30 and therefore exploratory.

At 10 ms this calibration CDF predicts 67.751\% misses, close to the observed
steady means of 67.569\% and 66.929\%. Across the twelve steady validation trials,
its MAE is 1.079 pp, versus 66.953 pp for frozen periodic+delay on that same
steady subset. Independently combining the worker and residual marginals gives
6.778 pp error; using median worker time with the residual distribution gives
9.822 pp. These contrasts show sensitivity to preserving the observed sum
distribution versus independently recombining stage marginals. They do not
establish a need for stage decomposition, or that the paired
CDF dominates at other deadlines: at exploratory 11 ms the independent
convolution is more accurate, and paired-CDF error rises to 4.169 pp at 10.5 ms.

The calibration CDF supplies no overload mechanism. At 25 ms it predicts zero
burst misses, producing 79.203 pp burst error. It is an ordinary reference for
the observed low-queue regime, not a universal alternative to periodic/FIFO
modeling. V31 below tests prospective accuracy within a declared low-queue
domain. An advance scope-selection rule for deployment still requires testing.

\section{Stage, accounting and mechanism boundaries}

At 10 ms the steady profiles produce 5310 misses. Computation is certainly
after the deadline for 2877, certainly within it despite a client miss for
1579, and clock-boundary ambiguous for 854. At 25 ms, all 6558 misses occur in
profile C: 6529 have computation certainly after the deadline, 22 certainly
within it despite a client miss, and seven are ambiguous. These are
descriptive counts, not a causal decomposition of CPU budget, frequency,
descheduling, dispatch or HTTP effects.

Treating the calibration worker-wall median as transferable CPU on-time while
retaining busy capacity would require 30.917 ms of worker availability per
100 ms, above the configured 25 ms CPU quota. The earlier rejected two-rate
identification therefore remains inadmissible. Neither clipping its budget nor
selecting faster calibration samples is justified. Aggregate CPU/job includes
the whole service; worker-wall time includes elapsed scheduling effects. A
measurement or intervention contract must distinguish these before a causal
claim is made.

\section{Prospective low-queue 10 ms confirmation}

V31 retains V30's service, quota, request work, binaries, controls and clock
joins, while independently repeating all sixteen calibration trials. Four
fresh low-rate calibration cohorts supply 676 successful queue<=100 us timing
samples (113/225/225/113). The queue-adjusted CDF (the artifact's
\path{joint_calibration_cdf}, also called paired CDF) uses $R=E-Q$;
all methods and per-offer probabilities across eleven declared deadlines are
sealed before the first validation trial. No V30 outcomes or historical
software-fixture parameters enter the V31 fit.
Its calibration input excludes observed queue time; its validation target
remains complete scheduled-to-client-completion time including waiting.

Eighteen trials use steady rates 5, 11 and 19 requests/s, six blocks in the
retained counterbalanced orders, 20 s warmup and 60 s measurement. There are
12600 measurement offers. The scientific domain is zero recorded measurement
throttled periods and at least 99\% of all offered requests with an observed
worker queue interval <=100 us. Domain exits must remain adverse results,
without replacement. These are operational limits checked on observations,
not proof of exclusive CPU service or an advance classifier of traffic state.

All eighteen trials are technically qualified and inside this domain. All
offers complete successfully; the smallest trial low-queue fraction is 99.825\%.
All 34 scientific trials complete without a technical retry. Full completion
replay and separate calculations from raw client/server records reproduce
calibration pairing, all-offer outcomes, domain fractions and block contrasts.

{\small\setlength{\tabcolsep}{2pt}\begin{longtable}{>{\raggedright\arraybackslash}p{0.25200\linewidth}>{\raggedright\arraybackslash}p{0.16200\linewidth}>{\raggedright\arraybackslash}p{0.16200\linewidth}>{\raggedright\arraybackslash}p{0.16200\linewidth}>{\raggedright\arraybackslash}p{0.16200\linewidth}}\toprule
Validation profile & Rate requests/s & Mean observed 10 ms miss \% & Paired-CDF forecast \% & Paired-CDF MAE pp \\
\midrule\endhead
A & 5 & 68.611 & 67.160 & 1.891 \\
B & 11 & 68.106 & 67.160 & 1.212 \\
C & 19 & 65.746 & 67.160 & 2.430 \\
\bottomrule\end{longtable}}

The primary equal-profile/block mean absolute error is 1.844 pp for the paired
CDF and 67.340 pp for periodic+delay. All six blocks improve. Mean paired
improvement is 65.495 pp, with the declared six-block 95\% t interval
[63.925,67.065] pp (df=5, multiplier 2.571). Every frozen confirmation gate
passes, including the <=3 pp mean block CDF error and >=5 pp lower improvement
endpoint. Maximum individual CDF trial error is 5.055 pp; the mean accuracy
gate is not an individual-trial error guarantee.

The CDF predicts the <=5\%/>5\% decision correctly in all eighteen trials at both
10 and 12 ms. All trials fail the miss budget at 10 ms; all have zero observed
misses at 12 ms. This confirms decisions well separated from the miss-budget
boundary. It does not establish precision near 5\% or a rate-selection method.
At these widely separated endpoints, constant labels of fail at 10 ms and
pass at 12 ms would give the same decisions. Decision agreement therefore
adds no evidence of locating an operating point near the miss-budget boundary.
At 12 ms the periodic forecast also has zero error; that endpoint adds no
comparative accuracy discrimination. Zero observed misses do not guarantee
zero future misses.

{\small\setlength{\tabcolsep}{2pt}\begin{longtable}{>{\raggedright\arraybackslash}p{0.25200\linewidth}>{\raggedright\arraybackslash}p{0.21600\linewidth}>{\raggedright\arraybackslash}p{0.21600\linewidth}>{\raggedright\arraybackslash}p{0.21600\linewidth}}\toprule
Deadline ms & Paired-CDF MAE pp & Independent worker/residual convolution MAE pp & Periodic+delay MAE pp \\
\midrule\endhead
9 & 1.940 & 1.180 & 69.519 \\
9.5 & 1.738 & 2.811 & 77.478 \\
10 (primary) & 1.844 & 7.585 & 67.340 \\
10.5 & 3.852 & 4.645 & 24.838 \\
11 & 0.607 & 0.628 & 2.243 \\
12 & 0.000 & 0.113 & 0.000 \\
\bottomrule\end{longtable}}

At the primary endpoint, median worker plus empirical response residual has
8.104 pp error; fluid busy and fluid quadratic each have 32.512 pp error.
The fixed-zero-phase/no-added-delay diagnostic has 67.488 pp error. These
secondary comparator contrasts are descriptive. The queue-adjusted CDF's advantage over the independent marginals at 10 ms
does not extend to every endpoint: the independent convolution is better at
9 ms, and paired-CDF error rises to 3.852 pp at 10.5 ms. The 3 pp primary gate
was never a gate across the deadline grid. Both CDF and periodic forecasts
have zero error at 12/20/25/50/100/500 ms in this study.

Cohort-only CDF 10 ms predictions range from 57.522\% to 69.912\%. Leaving each
calibration cohort out gives 1.974\textendash{}2.268 pp validation error. This sensitivity
is descriptive and does not alter the frozen pooled fit. Four calibration
trials and six validation blocks remain the experiment-level evidence; the
676 samples and 12600 offers are not additional independent replicates.

There are 8429 observed 10 ms misses: 4813 with computation certainly after
the deadline, 2264 certainly within despite the client miss, and 1352
clock-boundary ambiguous. At 11 ms, the corresponding counts are 279 misses,
12 certainly after, 233 certainly within and 34 ambiguous. These counts
conserve all misses and retain timing uncertainty, without causal attribution
to CPU demand, quota phase, descheduling or HTTP delay.

\subsection{Retrospective complete-response CDF control}

The manuscript review adds an ordinary control chosen after the V31 outcomes:
pool the complete client response times $E$ from the same four predeclared
timing-calibration cohorts and evaluate
$\hat p_E(D)=n^{-1}\sum_{i=1}^{n}\mathbf{1}\{E_i>D\}$ at the original frozen
10 ms deadline. All 676 calibration offers succeeded and met the original
queue<=100 us selection, so both pools contain exactly the same requests.
No cohort, threshold or validation duration is fitted or trimmed. Mean and
maximum calibration queue waits are 9.614 and 46.557 us, respectively.

{\small\setlength{\tabcolsep}{2pt}\begin{longtable}{>{\raggedright\arraybackslash}p{0.25200\linewidth}>{\raggedright\arraybackslash}p{0.21600\linewidth}>{\raggedright\arraybackslash}p{0.21600\linewidth}>{\raggedright\arraybackslash}p{0.21600\linewidth}}\toprule
Calibration reference at 10 ms & Forecast miss \% & Validation MAE pp & Maximum trial error pp \\
\midrule\endhead
Frozen queue-adjusted CDF, prospective & 67.160 & 1.844 & 5.055 \\
Complete-response CDF, retrospective & 67.456 & 1.802 & 5.350 \\
\bottomrule\end{longtable}}

At greater precision, the complete-response forecast is 67.455621\%, with MAE 1.802499 pp
and maximum trial error 5.350358 pp. Its mean error is only 0.041938 pp below
the frozen reference, while its maximum error is slightly larger. These
descriptive differences establish neither superiority nor statistical
equivalence. The narrow observed calibration queue waits accompany this
small difference; they do not establish a causal explanation.

This check limits any claim that queue subtraction or worker/residual pairing
is necessary for the good 10 ms forecast in these records. The original V31
comparison against periodic+delay remains prospective; this newly selected
control remains retrospective despite using only calibration samples for
its fit. No new confirmation gate is applied. All observations were collected
with the original instrumentation, and queue measurement remains part of
V31's domain qualification. Client-only deployment, removal of observers and
calibration-cost savings have not been qualified or prospectively tested.

\section{Subsequent NGINX application-transfer limit}

A separately sealed prospective study tests stock NGINX 1.30.5 dynamic gzip
compression of a fixed 1 MiB corpus on the same physical host, at gzip level 6
with native zlib 1.3.2. The application,
native runtime, single event-loop worker, CPU-2 affinity and 25 ms/100 ms quota
are pinned, with zero configured burst. Unlike V31, internal queue wait is
unobserved: the predictor uses complete scheduled-offer-to-body-receipt latency,
including all waiting. No queue subtraction or paired-stage periodic/residual
comparator is imported. This is a new application-transfer question, not an
exact replication of V31's model comparison.

Three fresh new-rate qualification cohorts precede four 400-offer calibration
cohorts at 5 requests/s. The lower pooled calibration median, 20,565,395 ns,
defines the primary deadline. The empirical CDF predicts 50\% misses at that
deadline. The calibration, deadline, CDF and forecasts are sealed before the
eighteen 300-offer held-out cohorts at 4, approximately 6 and 8 requests/s,
ordered in six balanced blocks. Qualification observations are excluded from
the scientific fit and validation. All 1,600 calibration and 5,400 validation
offers are correct; measurement throttle increments are zero in every cohort.
Dispatch, runtime identity, process/thread stability, request accounting,
observer inventory, owned cleanup and full raw replay pass. The campaign
completes without a technical or applicability exit in 29 minutes 2.660 seconds.

{\small\setlength{\tabcolsep}{2pt}\begin{longtable}{>{\raggedright\arraybackslash}p{0.25200\linewidth}>{\raggedright\arraybackslash}p{0.21600\linewidth}>{\raggedright\arraybackslash}p{0.21600\linewidth}>{\raggedright\arraybackslash}p{0.21600\linewidth}}\toprule
Validation rate requests/s & Mean observed primary miss \% & Frozen forecast \% & MAE pp \\
\midrule\endhead
4 & 48.833 & 50.000 & 1.389 \\
approximately 6 & 52.611 & 50.000 & 10.167 \\
8 & 43.611 & 50.000 & 6.389 \\
\bottomrule\end{longtable}}

Primary mean block absolute error is 5.981481 pp. The declared conditional,
model-based six-block 95\% t interval is [1.320988, 10.641975] pp; the required
upper endpoint is <=5 pp. The accuracy criterion is not supported. The interval
conditions on the finite frozen forecast and assumes approximately independent
normal block means. It does not propagate calibration uncertainty or establish
that unknown population mean error exceeds five points.

All six 8-requests/s cohorts have 41.000\textendash{}46.667\% misses, below the 50\% forecast.
The approximately-6-requests/s cohort in block 2 has 88.333\% misses, giving
38.333 pp error, and remains included. Its median full-response time is
21.438653 ms, with median dispatch lag 0.541647 ms and median dispatch-to-full-
body interval 20.857623 ms. That latter interval includes server waiting,
compression and delivery and is not an identified processing stage. Saved
records do not establish a scheduler, quota-phase or internal-queue cause.
The largest block contributes 41.176\% of the sum of the six block mean errors.
That concentration limits interpretation of the overall mean; it supplies
no justification for excluding the block or changing the primary criterion.

Fresh calibration cohort miss fractions at the pooled deadline are 44.500\%,
56.250\%, 48.500\% and 50.750\%; medians span 20.510439\textendash{}20.621183 ms. This
observed cohort variation limits a stationarity claim near the steep response
transition. Predeclared cohort-probability sensitivity MAEs are 5.296\textendash{}11.463 pp,
reported descriptively without changing the pooled forecast. At the fixed
secondary 10 ms deadline all validation offers miss; at 25 and 50 ms there
are zero observed misses. These secondary matches do not replace the failed
primary endpoint or guarantee zero future misses.

This completed experiment adds an adverse transfer result to the extension:
zero recorded measurement throttling and successful apparatus checks were
insufficient to ensure the declared accuracy of a rate-invariant complete-
response empirical forecast in this NGINX campaign. It does not invalidate
V31's confirmation within its different observed low-queue domain. Neither
new-algorithm novelty, causal mechanism nor independent-host reproduction is
established. All trials, the original primary deadline and the original gate
remain in the reported result.

A subsequent exploratory diagnosis uses only these sealed records. The
block-2 excursion's first 30 offers have 56.667\% misses; its remaining 270
have 91.852\%, showing persistence within the measurement window. Relative
to the other five same-rate cohorts, mean full response increases by 1.275873
ms: dispatch lag changes by -0.035438 ms, dispatch-to-headers by +0.189690 ms,
and headers-to-full-body by +1.121621 ms. About 87.9\% of the extra mean duration
therefore occurs after headers. These are observed client milestone intervals,
not identified server stages. Encoded response bytes/hashes are identical
across all 7,300 campaign measurement offers, and client response windows do
not overlap. Aggregate CPU per offer does not uniquely identify the excursion:
two other same-rate cohorts have comparable or higher average CPU use without
its large miss fraction.

At diagnostic deadlines 100 us below and above the frozen primary deadline,
the pooled calibration miss fractions are 57.750\% and 41.625\%. This steep
region makes small duration shifts consequential; these post-hoc offsets do
not replace the primary endpoint. Internal request-stage CPU/queue timings,
scheduler traces, client GC timing, frequency/thermal timelines and exact
quota phase were not recorded, so finer causal attribution remains unavailable.
The separate diagnosis report and per-request milestone CSV retain all offers
and explicitly distinguish these exploratory observations from confirmation.

\section{Fresh NGINX power-policy confirmation}

After the original NGINX transfer failure, a separately designed power-policy
pilot stopped at its zero-throttling applicability gate. Twelve of twenty-four
planned validation cohorts were collected; cohort 12 recorded throttling.
The entire prefix is preserved as exploratory, including the exit. It supplies
no complete six-block acceptance test and is not pooled with the following
fresh confirmation. This stop was not repaired by deleting or replacing a trial.

The fresh confirmation prospectively changes the population to all offered
requests under the configured quota, including throttling and waiting. Transport
errors and non-200 responses count as misses; corrupt complete output or invalid
identity/accounting records invalidate collection. The change is a disclosed
scope change, not a resume of the pilot. Pilot observations do not enter any
new reference. The preserved stock NGINX/client binaries, 1 MiB corpus, gzip
level 6, one worker on CPU 2 and 25 ms/100 ms quota remain fixed. The primary
deadline remains 20.565395 ms, rather than a new favorable endpoint.

P0 uses the powersave governor with balance\_performance EPP; P1 uses the same
governor with performance EPP. The setting is applied before ten-second settling
and held through each fresh service. Four new 400-offer calibration cohorts
at 5/s supply two references per policy in a randomly chosen balanced order.
All global forecasts freeze before validation. Six randomized blocks each
contain P0-4/s, P0-8/s, P1-4/s and P1-8/s once; each cohort has 25 warmup
and 300 measurement offers. Warmup and sample-weighted updates freeze separately
before that cohort's measurement and are later-decision comparisons.

Qualification and all twenty-four validation cohorts completed. The scientific
campaign retained 1,600 calibration, 7,200 validation and 700 warmup offers in
35.79 minutes. All responses succeeded; neither scientific phase recorded
throttling. Source/evidence seals, advance forecast barriers, all-offer accounting,
owned cleanup and restoration records verify. Absence of recorded throttling
does not remove the broader all-offer contract.

{\small\setlength{\tabcolsep}{2pt}\begin{longtable}{>{\raggedright\arraybackslash}p{0.25200\linewidth}>{\raggedright\arraybackslash}p{0.16200\linewidth}>{\raggedright\arraybackslash}p{0.16200\linewidth}>{\raggedright\arraybackslash}p{0.16200\linewidth}>{\raggedright\arraybackslash}p{0.16200\linewidth}}\toprule
Forecast & MAE pp & Conditional 95\% interval pp & Maximum error pp & Cohorts >5 pp \\
\midrule\endhead
Pooled policies & 24.729 & [20.377, 29.082] & 44.583 & 18/24 \\
Policy-matched & 15.708 & [12.695, 18.721] & 43.167 & 11/24 \\
P0-only & 55.458 & [52.445, 58.471] & 79.500 & 23/24 \\
25-offer warmup & 4.778 & [0.582, 8.974] & 25.000 & 7/24 \\
Matched + warmup & 15.250 & [12.286, 18.214] & 41.970 & 11/24 \\
Pooled + warmup & 24.350 & [20.036, 28.664] & 43.903 & 18/24 \\
\bottomrule\end{longtable}}

Policy matching improves over pooled calibration by 9.021 pp, conditional
95\% interval [1.746, 16.296] pp. The improvement condition passes, but the
matched accuracy upper bound 18.721 pp exceeds 5 pp: the combined claim is
not supported. Warmup's smaller mean error does not rescue the claim; its
upper bound is 8.974 pp, maximum error 25 pp, and information/budget differ.

P0 calibration predicts 79.5\% misses (81\% and 78\% in its two cohorts); later
P0 validation averages 48.889\%, ranging from 36.333\% to 84.333\%. P1 calibration
and all twelve P1 validation cohorts have zero observed primary misses. This
easy P1 endpoint does not hide P0's failure or certify zero future risk.
The pooled reference deliberately combines different assigned policies;
beating it is not superiority over a published state-aware predictor. The
matched references use 800 offers per policy, while both global procedures
share a 1,600-offer acquisition process. The baseline-only reference uses P0
samples for both policies and is another deliberately mismatched control.

The successful P0 mean response changes from 20.507 ms during calibration to
20.190 ms during validation, mainly through shorter dispatch-to-header elapsed
time; client stages do not identify server CPU or scheduling mechanisms.
Calibration is at 5/s and validation at 4/s or 8/s, leaving rate and chronology
unresolved. Matching EPP does not establish fixed physical frequency: the
\href{https://docs.kernel.org/admin-guide/pm/amd-pstate.html}{AMD active-driver documentation} \cite{amdpstate}
describes autonomous hardware decisions influenced by workload and physical
conditions. No historical frequency cause is established. At secondary 10 ms
all measured offers miss; at 25 and 50 ms all are timely. These saturated
endpoints do not replace the failed primary accuracy test.

\section{Fresh NGINX rate-transfer control under fixed P0}

A subsequent prospective control asks whether full target-rate calibration
transfers accurately and materially improves over ordinary 5/s calibration.
It retains the exact application, binaries, client endpoint, quota and
20.565395 ms deadline. P0 remains fixed with read-only policy checks; no host
power policy is written. Qualification includes three rate smoke cohorts and
a full 400-offer 4/s lifecycle case. Qualification and historical power outcomes
cannot supply scientific reference samples or choose a new endpoint.

Six fresh calibration cohorts contain two 400-offer cohorts each at 4, 5 and
8/s, in a random balanced order. All global references freeze before validation.
Six randomized blocks each contain one fresh 300-offer cohort at 4/s and one
at 8/s. Every cohort has 25 warmup offers; all-offer scoring, retained throttling,
response integrity and later warmup forecast barriers follow the frozen contract.
Qualification and all twelve validation cohorts completed. The scientific
campaign retained 2,400 calibration, 3,600 validation and 450 warmup offers in
23.96 minutes, all successful.

{\small\setlength{\tabcolsep}{2pt}\begin{longtable}{>{\raggedright\arraybackslash}p{0.25200\linewidth}>{\raggedright\arraybackslash}p{0.16200\linewidth}>{\raggedright\arraybackslash}p{0.16200\linewidth}>{\raggedright\arraybackslash}p{0.16200\linewidth}>{\raggedright\arraybackslash}p{0.16200\linewidth}}\toprule
Forecast & MAE pp & Conditional 95\% interval pp & Maximum error pp & Cohorts >5 pp \\
\midrule\endhead
Rate-matched & 1.194 & [0.808, 1.581] & 2.333 & 0/12 \\
Fixed 5/s & 1.972 & [1.037, 2.907] & 4.500 & 0/12 \\
Pooled rates & 1.361 & [1.091, 1.632] & 3.208 & 0/12 \\
25-offer warmup & 7.528 & [4.629, 10.426] & 13.667 & 7/12 \\
Matched + warmup & 1.149 & [0.841, 1.457] & 2.182 & 0/12 \\
5/s + warmup & 2.073 & [1.055, 3.091] & 4.485 & 0/12 \\
Pooled + warmup & 1.371 & [1.074, 1.669] & 3.216 & 0/12 \\
\bottomrule\end{longtable}}

Rate-matched accuracy passes the <=5 pp upper-bound condition. The improvement
over fixed 5/s is 0.778 pp, interval [-0.375, 1.931] pp, failing both the mean
>=2 pp and positive lower-bound advantage conditions. Ordinary 5/s and pooled
calibration also meet the accuracy target; every cohort error for all three
is <=5 pp. Neither superiority nor equivalence between rate-matched and 5/s
references is established. The combined accuracy-and-advantage claim is not
supported, while accurate rate-matched transfer in this episode is supported.

Each target-rate reference uses 800 calibration offers. Covering both 4/s and
8/s requires 1,600 offers, versus 800 for the shared 5/s reference; pooling
uses 2,400. All three pools were acquired by the actual study. Matched and
baseline effective sample counts are equal per target, but standalone procedural
costs were not experimentally compared. Warmup and updates decide later and
consume additional samples; their accuracy does not establish a free remedy.

Primary observed misses range from 78.667\% to 84.000\%. Accurate prediction
therefore describes poor deadline compliance, rather than useful SLO control
near the 5\% miss budget. The primary endpoint is not saturated. Validation
cohort 5 (8/s, block 3) records one throttled period and 6,600 us aggregate
throttling; cohort 9 (8/s, block 5) records one period and 9,133 us. Both remain
in all scores. Aggregate counters do not align an event with a specific response.
No other scientific observation flags or unsuccessful offers were recorded.
At secondary 10 ms every measured offer misses; at 50 ms all are timely. At
25 ms matched and 5/s MAE are 0.027778 pp. No secondary endpoint replaces
the predeclared primary result.

The earlier power P0 and the fresh 5/s P0 references both happen to predict
79.5\% misses from different fresh samples, but their later distributions differ.
Because ordinary 5/s calibration also transfers accurately here, these data do
not establish that matching offered rate repaired the preceding failure.
Campaigns differ in time, fresh instances and unmeasured physical state; they
are not randomized trials of the cause or repair. Keep the successful and
failed episodes separate, without pooling them into a favorable acceptance test.

\section{Fresh NGINX calibration-only deadline panel}

The earlier saved-data deadline-scope check showed why an accurate historical
endpoint may hide errors near the response CDF transition. That analysis was
retrospective development. The subsequent generation-01 study separately froze
five deadlines per policy from fresh calibration, before any validation outcomes.
It retained stock NGINX/native zlib, the 1 MiB gzip-6 corpus, Q25 ms/P100 ms,
one worker on CPU2 and the recorded P0/P1 policy contracts. Matching policy
readback does not hold physical execution speed fixed.

Four calibration services acquired 800 measured offers per policy. Zero-based
lower-order ranks 79, 199, 399, 599 and 719 froze the 1/10, 1/4, 1/2, 3/4 and
9/10 cutoffs. P0 cutoffs were 17.433213, 20.683338, 21.041565, 21.341755 and
21.562926 ms; P1 cutoffs were 16.133618, 16.383808, 16.664488, 16.946228 and
17.170689 ms. Six randomized blocks each included P0/P1 at 4/s and 8/s, with
25 warmup and 300 target offers per fresh service. The historical 20.565395 ms
cutoff was additional and descriptive; it did not select the primary result.

For each policy and block, average the absolute miss-probability errors across
both rates at each cutoff, then take the maximum across the five cutoffs.
The mean of six block scores is the primary statistic. Each policy's conditional
97.5\% interval must have upper endpoint <=5 pp. This max-after-rate-average
score differs from the historical single-deadline MAE and from averaging panel
errors. The paired policy family has >=95\% coverage only under independent,
approximately normal block scores. Six chronological blocks do not verify
those assumptions; shared calibration uncertainty is not propagated.

{\small\setlength{\tabcolsep}{2pt}\begin{longtable}{>{\raggedright\arraybackslash}p{0.25200\linewidth}>{\raggedright\arraybackslash}p{0.16200\linewidth}>{\raggedright\arraybackslash}p{0.16200\linewidth}>{\raggedright\arraybackslash}p{0.16200\linewidth}>{\raggedright\arraybackslash}p{0.16200\linewidth}}\toprule
Policy & Matched score (pp) & Conditional 97.5\% interval (pp) & Pooled score (pp) & Accuracy supported \\
\midrule\endhead
P0 & 6.611 & [1.831, 11.391] & 44.111 & No \\
P1 & 4.889 & [2.620, 7.158] & 38.222 & No \\
\bottomrule\end{longtable}}

The frozen status is \path{deadline_panel_accuracy_not_supported}. P1's point score
below 5 pp does not pass its upper-bound gate. Neither failed gate proves that
the population score exceeds 5 pp. Secondary pooled-minus-matched gains were
37.500 pp for P0 (conditional descriptive 95\% interval [33.703, 41.297]) and
33.333 pp for P1 ([30.405, 36.262]); they do not replace absolute accuracy.
Warmup-only scores were 14.278/11.556 pp, with upper bounds 20.186/17.228 pp.
Matched-plus-warmup scores were 6.515/4.886 pp, with upper bounds 11.261/7.166 pp.
These predeclared ordinary controls remain descriptive, without choosing a
winner after validation. Every control is charged the complete acquisition;
warmup updates decide later and use additional information.

Matched panel-average MAEs were 3.483/3.067 pp, masking harder individual cutoffs.
Each policy had 11 errors >5 pp among 60 cohort/cutoff cells; the maximum errors
were 19.667/11.000 pp. These 120 correlated cells are not independent replicates.
At the older cutoff P0 MAE was 3.181 pp and P1 was zero. P1 calibration and targets
both had zero misses there. Yet P1 cohort 15 at 4/s forecast 25\% versus 36\%
misses at its 16.946228 ms cutoff. This post-validation illustration shows the
scope of the predeclared panel result; it supplies no advance selection rule.

All 10,530 offers succeeded: 1,030 qualification, 1,600 calibration measurements,
100 calibration warmup, 7,200 target and 600 target warmup offers. Qualification
took 279.086 s and science 2,140.385 s, with source/forecast barriers, all-offer
accounting, owned cleanup and normal/dead-owner restoration verified. Validation
cohort 1 (P0, 8/s) recorded one throttled period / 4,703 us and remains included.
That aggregate observation cannot assign delay to requests or identify the
forecast error's cause. The complete result supports a deadline-scope limitation
of the tested reference, without a uniform CDF guarantee, identified mechanism,
new estimator or independent-host replication. Its frozen gates remain unchanged.

\section{Exploratory advance-reliability checks}

The original Go and NGINX saved records were examined for information available
before future measurement that could warn of poor transfer. These checks were
declared after the original outcomes were known and remain exploratory. Smaller
calibration dispersion does not rank the two episodes' later accuracy correctly:
Go's four complete-response calibration cohorts span 13.274 pp and NGINX's
span 11.750 pp, while their later MAEs are 1.802 and 5.981 pp respectively.
Ordinary cohort agreement, sampling-width illustrations and deadline sensitivity
checks either flag both episodes or fail to warn usefully. Same-population
sampling uncertainty does not bound unmeasured future distribution changes.

Recent 25-offer warmup CDFs give MAE 5.037 pp for Go and 9.463 pp for original
NGINX, worse than their older complete-response references. Ordinary pooled
updating modestly changes some Go errors but leaves the large NGINX excursion.
The fresh power and rate studies prospectively include their own ordinary
warmup comparisons; their respective warmup accuracy upper bounds of 8.974
and 10.426 pp also fail the 5 pp reference. Decision time and acquisition budget
must accompany these comparisons. No warning threshold, blend or warmup window
is selected from the validation outcomes.

Later saved-data screens retain their development status: the campaign outcomes
were known before the candidate checks were specified. Using information recorded
before target measurement does not make a newly chosen rule prospective.

{\small\setlength{\tabcolsep}{2pt}\begin{longtable}{>{\raggedright\arraybackslash}p{0.25200\linewidth}>{\raggedright\arraybackslash}p{0.32400\linewidth}>{\raggedright\arraybackslash}p{0.32400\linewidth}}\toprule
Fixed development screen & Retained observation & Limit \\
\midrule\endhead
Warmup CPU-duration shift & Power: three larger-error forecasts accepted; rate: all twelve accurate forecasts refused & Aggregate CPU agreement and an additive timing mapping do not assure subsequent transfer \\
Client-stage guard, 54 NGINX cohorts & Answers 1/18 original, 12/24 power and 2/12 rate; power answers are exclusively zero-miss P1 & No useful advantage over whole-response checks; inadequate transition coverage \\
Continuous first-100/next-200 probe & Power MAE falls 15.521 to 2.729 pp, but one agreeing probe has 31.5 pp future error & Later information can improve averages while internal agreement fails to assure persistence \\
Five-cutoff sensitivity screen, 78 NGINX cohorts & New panel median-shift guard accepts 10/24, including three larger errors, and refuses eight within-margin cohorts & No confirmed reliability signal; exploratory cohort-maximum labels differ from primary gates \\
\bottomrule\end{longtable}}

The continuous screen compares every method on the same final 200 offers,
so its original-reference MAE differs from the 300-offer primary. Power cohort
23 has 46\% and 50\% misses in two 50-offer probe halves, passing the fixed <=5 pp
agreement check. Its 100-offer forecast is 48\%, versus 79.5\% in the subsequent
200 offers. The case illustrates an observed persistence failure without
identifying physical drift. A 100-offer continuous probe consumes roughly
12.5\textendash{}25 seconds at 8\textendash{}4/s in addition to warmup, calibration and decision overhead.
Saved body completion timestamps support the offline split; live decision and
integrity-check timing were not qualified. P1's zero-miss cases lower the pooled
power error and do not validate P0 reliability.

In the new panel's sensitivity screen, cohort 1 has only -20.851 us warmup
median drift and 1.875 pp implied probability change, yet 19.667 pp future error.
Ordinary median correction lowers older power panel MAE from 18.514 to 4.930 pp
but worsens original NGINX from 5.437 to 7.200, rate from 2.617 to 6.466 and
new panel from 3.275 to 6.709 pp. These panel-average development scores do not
revise any frozen block gate. All cohorts, common references, chronology and
failure/throttle observations remain included; no threshold or favorable subset
is selected for confirmation.

No novel estimator or validated advance reliability signal emerges. This is
a limitation of the tested observations and checks, not a refutation of
state-conditioned prediction or conformal methods. Their additional inputs,
feedback and distributional assumptions must be honored in any later comparison.

\section{Measurement qualification and precision/cost limits}

The separate worker CPU/runnable-wait observer completed both boundary probes
and all 24 randomized overhead services: 7,890 successful offers, six paired
on/off blocks per rate. Identity, chronology, cgroup residual and quiet carryover
checks passed, while all six behavioral equivalence intervals exceeded their
frozen +/-2 pp probability or +/-100 us mean-time margins. The status remains
\path{observer_overhead_qualification_inconclusive}, with \path{qualified=false}.
All point estimates were inside their margins and every interval contained zero;
this establishes neither equivalence nor harmful observer effects. The observed
0.908\textendash{}1.462 ms operation time per warmup does not replace behavioral qualification.
This is a measurement study, separate from the six scientific campaigns;
worker-wait's predictive usefulness was never tested.

Offline warmup precision/cost analysis retained all 78 NGINX calibration/warmup
pairs without using target outcomes for its precision calculations. The new
panel's median full 95\% warmup-median interval span was 525.452 us, versus
empirical symmetric CDF sensitivity radii 47.708/43.496 us. Exact binomial
rank geometry assumes IID/stationarity, unverified here; location-based mapped
ranges are not future-error intervals and do not cover shape changes or later
drift. A conservative one-window IID DKWM illustration requires 738 offers
for a sufficient <=5 pp bound at 95\%, not a minimum sample size or power result.
Replacing 24 target warmups by 738 offers would add 17,112 offers and 53.475
nominal minutes. The conditional science projection grows from 35.673 to
89.148 minutes, beyond the 45-minute cap, and violates the 120-second enclosing
service limit at both rates even before startup/drain. Longer warmups remain held.

A later execution-rate candidate specifies worker-attributed APERF/MPERF ratio
and runtime measurements, exact identity/window boundaries and randomized
observer-on/off comparison. This is an offline design; no hardware counter
measurement or new observer implementation occurred. At most approximately eight
pairs per rate fit its budget under old worker-wait service-cost observations.
Substituting that different observer's historical variability into a 14-gate
simultaneous-interval design gives warmup halfwidths 268.804/262.564 us, versus
a proposed 20 us allowance. Even 100 us width-only geometry needs 32/31 pairs.
These are zero-centered planning scenarios, not actual new-observer variability,
power, necessary sample sizes or all-gate success probabilities. The proposed
\textasciitilde{}0.05436\% per-window clock precision goal additionally uses a hypothetical
20 ms CPU-service scale and an unverified additive clock-to-response mapping.
Counter counts alone do not certify that mapping or absolute window accuracy.

Hold further local hardware collection with these candidates. Reopening requires
an independently justified measurement/comparison redesign or an explicitly
changed budget, followed by a newly frozen qualification. Measurement qualification
would still require a separate demonstrated pre-validation-state-to-future-error
relation, an ordinary comparator with identical information and acquisition cost,
and fresh evaluation of a frozen rule. No scope-only collection, small repeat,
relaxed gate or fitted selector follows from the present consolidation. The
pending independent V31 host study retains its separate contract.

\section{Consolidated evidence and contribution boundary}

{\small\setlength{\tabcolsep}{2pt}\begin{longtable}{>{\raggedright\arraybackslash}p{0.25200\linewidth}>{\raggedright\arraybackslash}p{0.16200\linewidth}>{\raggedright\arraybackslash}p{0.16200\linewidth}>{\raggedright\arraybackslash}p{0.16200\linewidth}>{\raggedright\arraybackslash}p{0.16200\linewidth}}\toprule
Study & Deadline (ms) & Reference MAE (pp) & Comparator MAE (pp) & Frozen result \\
\midrule\endhead
Go V30 & 25 & Periodic + delay: 2.341 & Fluid busy FIFO: 73.599 & Prospective gates passed \\
Go V31 & 10 & Queue-adjusted CDF: 1.844 & Periodic + delay: 67.340 & Prospective gates passed in observed domain \\
NGINX original & 20.565395 & Complete-response CDF: 5.981 & No head-to-head comparator & Prospective accuracy gate not supported \\
NGINX power confirmation & 20.565395 & Policy-matched CDF: 15.708 & Pooled policies: 24.729 & Improvement passed; accuracy failed \\
NGINX rate transfer & 20.565395 & Rate-matched CDF: 1.194 & Fixed 5/s CDF: 1.972 & Accuracy passed; advantage not supported \\
\bottomrule\end{longtable}}

This is an evidence ledger, not a cross-application or general method ranking.
The five single-deadline campaigns above and the sixth panel campaign
(Section 14) use the same physical host; fresh calibration and service instances
are not new hosts. The panel uses a worst-cutoff block score, so its 6.611/4.889 pp
results must not be read as the single-deadline MAE in this ledger. Their comparator information, traffic, applicability
contracts and acceptance gates differ. The unusually large V30 contrast partly
reflects its fluid comparator's per-job floor of 28.640 ms, already above the
25 ms deadline. V31's simple retrospective response control limits the case
for richer decomposition. Policy matching's pooled comparator intentionally
mixes assignments, and rate matching supplies no required extra advantage over
an ordinary shared reference. These limits accompany the positive numbers.

The defensible contribution is quantitative prospective evidence of the tested
procedures' adequacy and calibration-transfer limitations, including unfavorable
controls. No first result, new estimator, universal boundary, identified physical
mechanism or published-method superiority follows. Seals, advance barriers and
retained failures strengthen credibility; reproducibility alone is not scientific
novelty. The breadth and originality of the empirical insight remain the main
submission risk, rather than the recorded arithmetic.

\section{Interpretation and remaining scope}

The evidence documents scoped differences in predictive adequacy. Average busy
capacity supports retained long-deadline forecasts; ordinary periodic service
repairs a confirmed 25 ms operating-condition contrast; near the 10 ms response
transition, the effective-demand approximation misses measured job-level
timing that a simple empirical calibration reference captures in discovery
and V31 prospectively confirms within its declared low-queue domain.
Traffic state matters alongside deadline resolution. This is not a theorem
that one particular ratio $D/P$ determines the correct model.

The retrospective complete-response control achieves nearly the same V31
mean error without subtracting calibration queue time. Thus the prospective
result supports the calibrated latency distribution against the tested
effective-demand approximation, without establishing necessary stage
decomposition or an instrumentation-saving method.

V31 establishes a scoped fresh confirmation; V30's empirical-CDF burst failure
remains negative evidence outside that scope. The observed domain checks do
not constitute an advance rule for choosing a model before traffic runs.
The subsequent stock-NGINX complete-response study did not meet its declared
accuracy criterion (Section 11). This supplies negative evidence for that newly
scoped application-transfer question; the broad client-completion outcome is
shared, but calibration inputs, completion checks, domain qualification and
comparators differ. The fresh policy control improves a pooled comparison
without sufficient accuracy; the fresh rate control gives accurate transfer
without the required extra advantage over 5/s calibration (Sections 12\textendash{}13).
The fresh panel fails both policy accuracy gates despite zero P1 error at the
older cutoff (Section 14). Later reliability screens and qualification/cost
audits leave the transfer difference unexplained and further local collection
held (Sections 15\textendash{}16).
Independent host transfer and qualified mechanism
measurements remain subsequent scope extensions. Further within-service sweeps
would not supply that missing hardware evidence. Consolidating both the scoped
confirmation and the transfer limit yields an empirical contribution candidate.
The targeted prior-work review finds the underlying methods
established; it does not certify originality of the precise empirical result.

\subsection{Limits of generalization}

The physical host is shared by every study. Neither counterbalanced ordering
nor CPU affinity supplies exclusive CPU execution or independent hardware
replication. Recorded power/thermal settling and process controls reduce some
variation but do not establish a fixed execution-cost distribution. Only the
specified rates, work sizes, deadlines and traffic schedules were tested.

The principal comparisons use different accuracy criteria and different
validation cohorts. V30 compares periodic service against constant-rate FIFO
at 25 ms; V31 compares a queue-adjusted CDF against periodic+delay at
10 ms; NGINX tests a full-response CDF without that comparator. Their errors
cannot be pooled into a general method ranking. The power confirmation compares
policy matching with a mixed reference, and rate transfer compares matched with
shared 5/s calibration, under broader all-offer contracts. The fresh five-cutoff panel has its own prospective primary contract. Older
secondary deadline grids, cohort sensitivities and saved-stage diagnoses are
descriptive; they supply
neither additional primary confirmation nor a causal explanation.

The targeted literature review identifies close precedents for periodic service,
profiling and temporal dependence. No published competing implementation was
executed under a matched contract. The present evidence supports an incremental
empirical contribution candidate; it does not establish a first result or
superiority over published methods.

The user-arranged independent V31 replication is pending; no new-host outcome
or inspected completed package has been received. It is excluded from every
table and conclusion here. Further contribution work needs a concrete additional
capability, a compatible comparator with fair information/acquisition costs,
and an independent evaluation design. A causal extension additionally requires
measurements or interventions that identify the proposed cause. The current
result supplies no validated advance observable, so a new guard or another
same-settings campaign is not selected as a demonstrated contribution path.

\section{Reproducibility}

Frozen protocols, source/binary identities, per-offer forecasts, calibration
records, raw all-offer outcomes and result seals are retained. Completed
independent audits replayed the campaigns and checked raw timing/accounting.
The final manuscript audit reconciles tables, intervals and figures with those
unchanged records; it is not another experiment or independent-host replication.

The supporting material maps each
claim to its evidence and retains the failed application qualifications. The
reproduction guide supplies offline
verification commands and distinguishes manuscript checks from full campaign
replay. Five numerical checks of empirical boundary counting, convolution,
calibration-leakage rejection and stage-count conservation complement the
periodic slot-oracle tests. Software checks are not empirical replicates.

The 11 October consolidation adds the fresh deadline-panel result, later
development screens and measurement/precision holds to the complete power and
rate controls. A new offline audit verifies the new primary table and
Markdown/LaTeX derivation, replays the panel and qualification records and
preserves all historical tables and 2,559 protected inputs. Earlier checkers
run against their captured manuscript versions in a temporary logical workspace;
their original source and seals remain unchanged. Compiler diagnostics are recorded for the current source.
No new performance collection or public release accompanies this revision.

\section{Conclusion}

Under the tested fractional CPU quota, predictive adequacy depends on the
selected deadline and observed traffic conditions. A periodic-budget
approximation gives a confirmed 25 ms improvement over constant-rate FIFO,
while a separately calibrated empirical timing reference gives a confirmed
10 ms improvement within the Go service's declared low-queue domain. The
empirical reference has no overload mechanism and fails on the retained burst
traffic. A separately scoped stock-NGINX complete-response forecast does not
meet its declared primary accuracy criterion despite successful measurement
gates and zero recorded measurement throttling. The retrospective Go
complete-response control further limits claims that stage decomposition or
queue subtraction is needed for the observed 10 ms accuracy.

Fresh NGINX follow-ups preserve the transfer limitation: assigned-policy matching
improves the pooled reference but fails absolute accuracy, while later rate-matched
calibration and ordinary 5/s calibration both transfer accurately without a
supported incremental advantage. These outcomes do not identify a causal repair
or an advance rule for trustworthy transfer. The fresh panel further fails
both policy accuracy gates despite accurate prediction at an older saturated
cutoff. Later screens, observer qualification and precision/cost audits supply
no affordable qualified reliability method; further local hardware collection
with these candidates stays held. Independent-host replication remains
pending and is not included.

Together these results document the usefulness and limits of selected
established approximations in reproducible prospective tests. They support
neither a universal model-selection rule nor a quota-phase or execution-cost
mechanism. Further mechanism, application or host studies would address new
questions; the completed comparisons retain their original outcomes and scope.

\begin{thebibliography}{99}
\bibitem{S01} Linux kernel documentation. CFS Bandwidth Control. Live documentation accessed 2026-10-10. \url{https://docs.kernel.org/scheduler/sched-bwc.html}.
\bibitem{S02} Haoran Li; Chenyang Lu; Christopher Gill. Predicting Latency Distributions of Aperiodic Time-Critical Services. RTSS 2019. \url{https://www.cse.wustl.edu/~cdgill/publications/rtss19.pdf}.
\bibitem{S03} Tommaso Cucinotta; Luca Abeni; Mauro Marinoni; Riccardo Mancini; Carlo Vitucci. Strong Temporal Isolation Among Containers in OpenStack for NFV Services. IEEE Transactions on Cloud Computing, 2023. \url{https://retis.santannapisa.it/people/nino/publication/TCC23.pdf}.
\bibitem{S04} Ahsan Ali; Riccardo Pinciroli; Feng Yan; Evgenia Smirni. CEDULE: A Scheduling Framework for Burstable Performance in Cloud Computing. ICAC 2018. \url{https://www2.cs.uh.edu/~fyan/Paper/Feng-ICAC18.pdf}.
\bibitem{S05} Ningfang Mi; Giuliano Casale; Ludmila Cherkasova; Evgenia Smirni. Sizing multi-tier systems with temporal dependence: benchmarks and analytic models. Journal of Internet Services and Applications, 2010. \url{https://link.springer.com/content/pdf/10.1007/s13174-010-0012-9.pdf}.
\bibitem{S06} Andrew Jeffery; Chris Jensen; Richard Mortier. Reducing Tail Latencies Through Environment- and Neighbour-aware Thread Management. arXiv:2407.11582v1, 2024. \url{https://arxiv.org/html/2407.11582v1}.
\bibitem{capacityprecursor} Capacity Calibration and Deadline Forecasts Under a Fractional CPU Quota. Retained precursor manuscript, 2026; authors and public artifact location unset. Local artifact: \path{paper/capacity_calibration_empirical.md}.
\bibitem{tales} Jialin Li; Naveen Kr. Sharma; Dan R. K. Ports; Steven D. Gribble. Tales of the Tail: Hardware, OS, and Application-level Sources of Tail Latency. SoCC 2014. \url{https://syslab.cs.washington.edu/papers/latency-socc14.pdf}.
\bibitem{retail} Shuang Chen; Angela Jin; Christina Delimitrou; Jose F. Martinez. ReTail: Opting for Learning Simplicity to Enable QoS-Aware Power Management in the Cloud. HPCA 2022. \url{https://people.csail.mit.edu/delimitrou/papers/2022.hpca.retail.pdf}.
\bibitem{mpora} Abigail Eisenklam et al. Uncertainty-Aware Resource Allocation for Multi-Path Programs with In-Kernel Predictions. ECRTS 2026. DOI: 10.4230/LIPIcs.ECRTS.2026.17. \url{https://drops.dagstuhl.de/storage/00lipics/lipics-vol375-ecrts2026/html/LIPIcs.ECRTS.2026.17/LIPIcs.ECRTS.2026.17.html}.
\bibitem{amdpstate} Linux kernel documentation. amd-pstate CPU Performance Scaling Driver. Live documentation accessed 2026-10-10. \url{https://docs.kernel.org/admin-guide/pm/amd-pstate.html}.
\end{thebibliography}
\end{document}
